\documentclass[10pt]{article}

% Specify the margins and text size.
\setlength{\textwidth}{6.4in}
\setlength{\textheight}{9.5in}
\setlength{\oddsidemargin}{0pt}
\setlength{\evensidemargin}{0pt}
\setlength{\topmargin}{0pt}
\setlength{\hoffset}{.05in}
\setlength{\voffset}{-.8in}

\setlength{\parskip}{5pt}
\setlength{\parindent}{0pt}

% Load some fonts and symbol packages
\usepackage{latexsym}
\usepackage{pifont}       % contains 'star' symbol for counterinsurgency handbook title
\usepackage{yfonts} 
\usepackage{amsmath}
\usepackage{amsfonts}

\usepackage{graphicx}     % actually, this is loaded by pstricks
\usepackage[T1]{fontenc}
\usepackage{ifthen}
\usepackage{pstricks,pst-func,pst-text,pst-node,multido,pst-plot,calc,pst-3dplot}
%\usepackage[all]{xy}
%\usepackage{animate}

% The hyperref package inserts links into the pdf.
\definecolor{MyLinkColor}{rgb}{.1,.2,1}
\definecolor{MyCiteColor}{rgb}{.1,1,.2}
\definecolor{MyURLColor}{rgb}{.4,.4,.4}
\usepackage[backref=true,pagebackref=false,hyperindex,colorlinks=true,
  linkcolor=MyLinkColor,urlcolor=MyURLColor]{hyperref}


% The tweaklist package is something I found on the web.  It provides a simple interface
% for making changes to spacing used in the itemize and enumerate environments.  Comment
% this out if you don't care to use tweaklist.
\usepackage{tweaklist}
\renewcommand{\itemhook}{\setlength{\parskip}{2pt}\setlength{\parsep}%
{1pt}\setlength{\topsep}{0pt}\setlength{\itemsep}{0pt}}

\newcommand{\N}{\ensuremath\mathbb{N}}
\newcommand{\Q}{\ensuremath\mathbb{Q}}
\newcommand{\R}{\ensuremath\mathbb{R}}
\newcommand{\C}{\ensuremath\mathbb{C}}
\newcommand{\z}{\hphantom{0}}

\begin{document}
\pagestyle{empty}

\href{http://www.su.edu}{\includegraphics[height=1.75cm]{sulogo.eps}}
\vspace{-1.69cm}

{\small
\begin{tabular}{cl}
\hspace{5in} & Math 207:  Statistics\\
                & Warmup Exercises\\
\end{tabular}
}

\setlength{\baselineskip}{1.05\baselineskip}
\bigskip

1.  Without using a calculator, 
say whether each of the following is about 1\%, 10\%, 25\% or 50\% ---
\begin{center}
\begin{tabular}{clccl}
(a) & 39 out of 398 & \hspace{2in}  & (b) 99 out of 407\\[10pt]
(c) & 57 out of 209 &               & (d) 99 out of 197
\end{tabular}
\end{center}

2. Among beginning statistics students in one university, 46 students
out of 446 reported family incomes ranging from \$40,000 to \$50,000 a year.
Answer the following without a calculator.
\begin{enumerate}
\item[(a)] About what percentage had family incomes in the range \$40,000 
  to \$50,000?
\item[(b)] Guess the percentate that had family incomes in the range \$45,000 to \$46,000
   a year.
\item[(c)] Guess the percentage that had family incomes in the range \$46,000 to
   \$47,000 a year.
\item[(d)] Guess the percentage that had family incomes in the range \$47,000 to
   \$49,000 a year.
\end{enumerate}

\newcommand{\HS}{\hspace{30pt}}
3. Draw a pair of axes and plot each of the following points:
\[(1,1)\HS (2,2)\HS (3,3)\HS (4,4)\]
What can you say about the points?

%%%%%%%%%%%%%%%%%%%%%%%%%%%%%%%%%%%%%%%%%%%%%%%%%%%%%%%%%%%%%%%%%%%%%%%%%%%%

\vfill

\href{http://www.su.edu}{\includegraphics[height=1.75cm]{sulogo.eps}}
\vspace{-1.69cm}

{\small
\begin{tabular}{cl}
\hspace{5in} & Math 207:  Statistics\\
                & Warmup Exercises\\
\end{tabular}
}

\setlength{\baselineskip}{1.05\baselineskip}
\bigskip

1.  Without using a calculator, 
say whether each of the following is about 1\%, 10\%, 25\% or 50\% ---
\begin{center}
\begin{tabular}{clccl}
(a) & 39 out of 398 & \hspace{2in}  & (b) 99 out of 407\\[10pt]
(c) & 57 out of 209 &               & (d) 99 out of 197
\end{tabular}
\end{center}

2. Among beginning statistics students in one university, 46 students
out of 446 reported family incomes ranging from \$40,000 to \$50,000 a year.
Answer the following without a calculator.
\begin{enumerate}
\item[(a)] About what percentage had family incomes in the range \$40,000 
  to \$50,000?
\item[(b)] Guess the percentate that had family incomes in the range \$45,000 to \$46,000
   a year.
\item[(c)] Guess the percentage that had family incomes in the range \$46,000 to
   \$47,000 a year.
\item[(d)] Guess the percentage that had family incomes in the range \$47,000 to
   \$49,000 a year.
\end{enumerate}

3. Draw a pair of axes and plot each of the following points:
\[(1,1)\HS (2,2)\HS (3,3)\HS (4,4)\]
What can you say about the points?

\eject
{\ }

Exercise: Simplify each of the following:
\newcommand{\MHS}{\hbox{\hspace{1.75in}}}
\begin{center}
\begin{tabular}{cccc}
a) $\frac{1}{3} + \frac{1}{2}$\MHS & b) $\frac{1}{4}+\frac{1}{2}+\frac{1}{5}$\MHS 
  & c) $\frac{1/2}{2}$\\[2in]
d) $\displaystyle\sqrt{\frac{4}{9}}$\MHS        & e) $\displaystyle\frac{1/2 + 1/5}{3}$\MHS   & 
   f) $\displaystyle\sqrt{1-\left(\sqrt{3}/2\right)^2}$
\end{tabular}
\end{center}


\vfill

Exercise: Simplify each of the following:

\begin{center}
\begin{tabular}{cccc}
a) $\frac{1}{3} + \frac{1}{2}$\MHS & b) $\frac{1}{4}+\frac{1}{2}+\frac{1}{5}$\MHS 
  & c) $\frac{1/2}{2}$\\[2in]
d) $\displaystyle\sqrt{\frac{4}{9}}$\MHS        & e) $\displaystyle\frac{1/2 + 1/5}{3}$\MHS   & 
   f) $\displaystyle\sqrt{1-\left(\sqrt{3}/2\right)^2}$
\end{tabular}
\end{center}
\vspace{1.25in}

\eject

%%%%%%%%%%%%%%%%%%%% VERSION 2

\href{http://www.su.edu}{\includegraphics[height=1.75cm]{sulogo.eps}}
\vspace{-1.69cm}

{\small
\begin{tabular}{cl}
\hspace{5in} & Math 207:  Statistics\\
                & Warmup Exercises\\
%                & 11 January 2012
\end{tabular}
}

\setlength{\baselineskip}{1.05\baselineskip}
\bigskip\bigskip

{\color{blue} Solutions are in blue.}
\medskip

1.  Without using a calculator, 
say whether each of the following is about 1\%, 10\%, 25\% or 50\% ---
\begin{center}
\begin{tabular}{clccl}
(a) & 39 out of 398 & \hspace{2in}  & (b) 99 out of 407\\[10pt]
(c) & 57 out of 209 &               & (d) 99 out of 197
\end{tabular}
\end{center}

{\color{blue}
%\textbf{Solutions:}

(a) $\frac{39}{398}\approx \frac{40}{400}=\frac{4}{40} =\frac{1}{10} = 10\%$.

(b) $\frac{99}{407}\approx \frac{100}{400}=\frac{1}{4}=25\%$.

(c) $\frac{57}{209}\approx \frac{50}{200}=\frac{5}{20}=\frac{1}{4}=25\%$.

(d) $\frac{99}{197}\approx \frac{100}{200}=\frac{1}{2}=50\%$.
}
\bigskip


2. Among beginning statistics students in one university, 46 students
out of 446 reported family incomes ranging from \$40,000 to \$50,000 a year.
Answer the following without a calculator.
\begin{enumerate}
\item[(a)] About what percentage had family incomes in the range \$40,000 
  to \$50,000?\\
 {\color{blue} $\frac{46}{446}\approx \frac{46}{460}=\frac{1}{10}=10\%$.}
\item[(b)] Guess the percentate that had family incomes in the range \$45,000 to \$46,000
   a year.\\
  {\color{blue} The total range from (a) was \$45,000 -- \$50,000.  So, the range for (b)
    is $1/10$ of that range.  If we take $1/10$ of the 10\% total, then we get 1\%.
    If we think about the boxes of a histogram.  the \$45K -- \$50K box has width 10 and area 10\%.
    So it has height 1 ($w\times h = a$). The slice between \$45K and \$46K has width 1 and height 1.}
\item[(c)] Guess the percentage that had family incomes in the range \$46,000 to
   \$47,000 a year.\\
  {\color{blue}  The range \$46,000 -- \$47,000 is $1/10$ of the total again.  So, its area is also 1\%.}
\item[(d)] Guess the percentage that had family incomes in the range \$47,000 to
   \$49,000 a year.\\
{\color{blue} This interval has width 2 so the area is 2\%.}
\end{enumerate}

%\newcommand{\HS}{\hspace{30pt}}
3. Draw a pair of axes and plot each of the following points:\vspace{-3pt}
\[(1,1)\HS (2,2)\HS (3,3)\HS (4,4)\vspace{-5pt}\]
What can you say about the points?

\bigskip
\begin{center}
{\color{blue}\begin{pspicture}(0,0)(4,3)
\psset{yunit=0.75,linecolor=blue}
\psaxes(0,0)(4,0)(0,4)
\psdot(1,1)\rput[l](1.2,1){$(1,1)$}
\psdot(2,2)\rput[l](2.2,2){$(2,2)$}
\psdot(3,3)\rput[l](3.2,3){$(3,3)$}
\psdot(4,4)\rput[l](4.2,4){$(4,4)$}
\rput[l](5,2){The points are all on a line.}
\end{pspicture}}
\end{center}

\vfill
\eject

Exercise: Simplify each of the following:

\renewcommand{\MHS}{\hphantom{\hspace{1.in}}}
\begin{center}
\begin{tabular}{llll}
a) $\frac{1}{3} + \frac{1}{2}$\MHS & b) $\frac{1}{4}+\frac{1}{2}+\frac{1}{5}$\MHS 
  & c) $\frac{1/2}{2}$\\[10pt]
{\color{blue} a) $\frac{1}{3} + \frac{1}{2} = \frac{2}{6} + \frac{3}{6} = \frac{5}{6}$} &
{\color{blue} b) $\frac{1}{4}+\frac{1}{2} + \frac{1}{5}=\frac{5}{20}+\frac{10}{20}+\frac{4}{20} = \frac{19}{20}$}
&
{\color{blue} c) $\frac{1/2}{2}=\frac{1/2}{2}\times \frac{2}{2} = \frac{(1/2)\times 2}{2\times 2}=\frac{1}{4}$}
\\[1in]
d) $\displaystyle\sqrt{\frac{4}{9}}$\MHS        & e) $\displaystyle\frac{1/2 + 1/5}{3}$\MHS   & 
   f) $\displaystyle\sqrt{1-\left(\sqrt{3}/2\right)^2}$\\[10pt]
{\color{blue} d) $\sqrt{\frac{4}{9}} = \frac{\sqrt{4}}{\sqrt{9}} = \frac{2}{3}$} &
{\color{blue} e) $\displaystyle\frac{1/2 + 1/5}{3} = \frac{\frac{5}{10} + \frac{2}{10}}{3}
    = \frac{7/10}{3} = \frac{7/10}{3}\times \frac{10}{10} = \frac{7}{30}$}
\end{tabular}
\end{center}

{\color{blue} f) $\displaystyle\sqrt{1 - \left(\sqrt{3}/2\right)^2} = 
    \sqrt{1 - \left(\frac{\sqrt{3}}{2}\right)^2} = 
    \sqrt{1 - \left(\frac{3}{4}\right)} = 
    \sqrt{\frac{1}{4}} = \frac{1}{2}$.  Notice that to do $\left(\frac{\sqrt{3}}{2}\right)^2$, 
    we have to square the top and the bottom.  So, $(\sqrt{3})^2 = 3$ and $(2)^2 = 4$.
\vspace{1.25in}


\end{document}
